\section{Bounded control: Booleanization, traces, search, and descent}

\subsection{Characteristic terms for fixed quantifier-free formulas}

A fixed quantifier-free formula over decidable arithmetic can be compiled into
a primitive-recursive characteristic term.  This fact is repeatedly used to
replace bounded logical control by equations.

\begin{theorem}[Quantifier-free Booleanization]\label{thm:booleanization}
For every fixed quantifier-free formula \(\theta(\vec x)\) whose atomic
predicates are equations or decidable relations represented by PR
characteristic functions, there is a PR term \(\chi_\theta(\vec x)\in\{0,1\}\)
such that iPRA derives
\[
   \chi_\theta(\vec x)=1\ \leftrightarrow\ \theta(\vec x)
   \qquad\text{and}\qquad
   \chi_\theta(\vec x)=0\ \leftrightarrow\ \neg\theta(\vec x).
\]
\end{theorem}

\begin{proof}
Proceed by structural induction on the fixed formula.  For equality, use the
PR equality test on natural-number codes and prove its defining equations by
quantifier-free induction.  Negation is Boolean complement, conjunction is
multiplication, and disjunction may be represented by truncated Boolean sum or
\(\max\).  Implication between decidable formulas is
\((1-a)\vee b\).  Bounded quantifiers, when present as abbreviations, are
compiled by bounded products or bounded sums using primitive recursion.  Each
correctness proof is quantifier-free induction on the finite bound.  The
formula is fixed, so this is syntax-directed compilation, not an internal
truth predicate for arbitrary formulas.
\end{proof}

The existing \path{ConciseIPRA} reflection files provide related
characteristic-term infrastructure.  The theorem above specifies the interface
needed by the graph compiler: every branch condition and bounded verifier is
reduced to an equation of PR terms.

\subsection{Clocked traces}

Let \(S:\N\times\N\to\N\) be a PR step function, where the first input is the
clock and the second is a coded state.  Let
\[
  s(0)=s_0,\qquad s(t+1)=S(t,s(t)).
\]
Because \(S\) is primitive recursive, the trace function \(s(t)\) is primitive
recursive.  It may be expressed directly by a PR symbol in ordinary iPRA or by
the native iterator after pairing the time and state when the step depends on
both.

\begin{theorem}[Finite trace reflection]\label{thm:trace}
Let \(J(t,s)\) be a fixed quantifier-free invariant.  Suppose iPRA derives
\[
   J(0,s_0)
   \quad\text{and}\quad
   J(t,s(t))\to J(t+1,S(t,s(t))).
\]
Then iPRA derives \(J(T,s(T))\) for every term \(T\), and derives the universal
closure \(\forall t\,J(t,s(t))\) when the free-variable form is used.
\end{theorem}

\begin{proof}
Apply quantifier-free induction to the formula \(J(t,s(t))\).  The base is the
first hypothesis.  The step uses the defining equation for the trace and the
second hypothesis.  If only the terminal instance is required, instantiate the
open induction theorem at \(T\).  In the native backend the same closure is
obtained via the introduction rule and the existing translation from ordinary
quantifier-free induction to \(\KPrf\).
\end{proof}

A trace node stores not only the final state but also the proof-relevant fields
needed by \(J\): flags that every array access was in bounds, accumulated error
budgets, query logs, and a certificate status.  This ``single state, single
invariant'' form avoids a separate induction for each component.

\subsection{Least bounded search}

Given a PR predicate \(P(x)\in\{0,1\}\) and a bound \(B\), define a state
\((f,w)\), where \(f\) records whether a witness has been found and \(w\) stores
the first one.  At time \(t\), if \(f=0\) and \(P(t)=1\), set
\((f,w)=(1,t)\); otherwise retain the state.  Iterate through \(B\).

\begin{theorem}[Least bounded search]\label{thm:bounded-search}
There is a primitive-recursive function \(\mu_{\le B}P\) with a status flag such
that iPRA proves:
\begin{enumerate}
\item if the flag is one, then \(w\le B\), \(P(w)=1\), and
\(P(v)=0\) for every \(v<w\);
\item if the flag is zero, then \(P(v)=0\) for every \(v\le B\);
\item if \(\exists v\le B\,P(v)=1\), the flag is one.
\end{enumerate}
\end{theorem}

\begin{proof}
Use the clocked search above.  The quantifier-free invariant says: (i) the flag
is Boolean; (ii) if set, the stored value is less than the clock, satisfies
\(P\), and all earlier values fail; (iii) if unset, all values below the clock
fail.  Bounded universal assertions in the invariant are Booleanized by finite
products.  The step is a finite case split on two Boolean equations.  At clock
\(B+1\), the invariant yields the three conclusions.
\end{proof}

The theorem supplies a selector only because the candidate space is explicitly
finite.  It is not countable choice and does not select from an unbounded
sequence of arbitrary inhabited sets.

\subsection{Complete bounded enumeration}

A similar trace builds a bit vector or coded list containing exactly the
witnesses \(v\le B\) with \(P(v)=1\).  The invariant states that the list is
strictly increasing, all its entries satisfy \(P\), and every satisfying value
below the clock occurs in the list.

\begin{theorem}[Complete enumeration]\label{thm:enumeration}
For every bounded decidable predicate, there is a PR list constructor and an
iPRA proof that its output is sound, duplicate-free, ordered, and complete
within the bound.
\end{theorem}

Complete enumeration is often preferable to least selection.  It preserves all
branches for later invariant observation and makes the source of finiteness
visible.

\subsection{Finite-state dynamic programming}

Many bounded algorithms are naturally expressed as dynamic programs: Gaussian
elimination over a fixed matrix size, coefficient recurrences, finite graph
reachability, subset recursions, or backward viability.  A table of length
\(M\) is coded by nested pairing or a beta-function representation.  The update
operation is PR because every lookup and replacement is bounded.

The general proof pattern is:
\begin{enumerate}
\item define a table invariant asserting correctness for all indices below the
current clock;
\item Booleanize that bounded assertion;
\item prove the next cell correct from previously established cells;
\item apply Theorem~\ref{thm:trace}.
\end{enumerate}
This is more robust than attempting to reason about a host array and later
reflecting its output.

\subsection{Lexicographic descent and explicit fuel}

Algorithms in algebra and number theory frequently recurse by descent on a
fixed tuple of natural measures: degree, valuation, denominator, support size,
or a pair of dimensions.  A strict lexicographic descent in a bounded box has an
explicit finite fuel.

\begin{theorem}[Finite lexicographic fuel]\label{thm:lexfuel}
Let
\[
  L=\prod_{i<k}\{0,\ldots,B_i\}
\]
with fixed \(k\).  Every strictly descending sequence in lexicographic order
has length at most
\[
       |L|=\prod_{i<k}(B_i+1).
\]
Consequently a descent algorithm on \(L\) can be compiled to a clocked trace
with that fuel.
\end{theorem}

\begin{proof}
Lexicographic order is a strict total order on the finite set \(L\).  A strictly
descending sequence has no repeated state.  There are \(|L|\) states, so its
length is at most \(|L|\).  Constructively, map a tuple to its mixed-radix rank
\[
 r(x_0,\ldots,x_{k-1})
   =x_0\prod_{j>0}(B_j+1)+x_1\prod_{j>1}(B_j+1)+\cdots+x_{k-1}.
\]
Lexicographic decrease implies natural-number rank decrease.  The rank formula
and its monotonicity proof are primitive recursive and quantifier-free for fixed
\(k\).
\end{proof}

If a classical proof says merely that a well-founded measure decreases, the
iPRA reconstruction must expose a finite bound on the relevant initial segment.
The mixed-radix construction provides one for fixed-dimensional tuples.

\subsection{Finite exceptional-set synchronization}

Several estimates may each fail at finitely many indices.  One often needs a
single index where all estimates hold.

\begin{theorem}[Common good finite window]\label{thm:common-window}
Let \(E_0,\ldots,E_{r-1}\subseteq\N\) be finite sets with
\(|E_i|\le b_i\).  For every \(N\), the interval
\[
       [N,N+\textstyle\sum_i b_i]
\]
contains an index outside every \(E_i\).
\end{theorem}

\begin{proof}
The interval has \(1+\sum_i b_i\) elements.  The union of the exceptional sets
contains at most \(\sum_i b_i\) elements, because a coded injection of each
\(E_i\) into a block of size \(b_i\) gives a coded surjection from a set of that
size onto the union.  Therefore at least one interval point is not in the
union.  Bounded least search finds the first.  All counting is over coded finite
lists and is provable by quantifier-free induction.
\end{proof}

This theorem resolves a common ``choose a sufficiently large common index''
gap when explicit exception budgets are available.  Without such budgets,
finite exceptionality alone does not provide a computable search window.

\subsection{Metastability rather than an unproved modulus}

A sequence may be known classically to converge without a primitive-recursive
modulus available to the proof.  A valid finite substitute is metastability:
for every \(q\) and every control function \(g\), produce \(N\le B(q,g)\) such
that
\[
    \forall i,j\in[N,N+g(N)]\quad d(a_i,a_j)\le2^{-q}.
\]
When \(B\) is explicit and PR, the finite window can be searched and checked.
A genuine modulus is stronger and should not be inferred from metastability
without an additional theorem.
